\section{Methods}

The inspection crawler's control software processes each cycle in four steps. First,
the raw wall-distance signal from the two ToF sensors is read at 100 Hz. Second, the
signal is low-pass filtered using a single-pole IIR filter with alpha=0.2. Third, the
filtered distance is passed to the wall_follow controller (ctrl_mode="wall_follow" in
the config), which computes a lateral correction using a PD law with kp=2.1 and
kd=0.35; the correction is then clipped to 0.15 m/s. Fourth, the track velocities are
sent to the motor driver via the set_track_velocity() API of the vendor SDK at 100 Hz.

The state estimator is an extended Kalman filter with process noise
Q=diag(0.01, 0.01, 0.001), measurement noise R=diag(4.0, 4.0), and initial covariance
P0=100*I. The EKF update is capped at 40 iterations, and it's checked for divergence
every 5 cycles. A simulation mode shares the same controller code but was not used for
the results reported here. The controller is inspired by impedance control but is not
formally an impedance controller, because the contact dynamics are not modeled. The
selection of the filter parameters is described in Section IV.

\section{Results}

We conducted sixteen runs through the 12 m test pipe. In each run, the crawler was
placed at the pipe entrance, the operator started the wall_follow controller, and the
crawler traversed the pipe autonomously at 0.1 m/s while logging wall distance at 100
Hz; runs alternated between the compensated and uncompensated controller.

Figure 3 shows the wall-distance error over time. The compensated controller kept the
mean absolute error at 3.1 mm, which shows that the thermal compensation is effective
and confirms our design choice of placing the reference sensor at the front bulkhead.
By direction, clockwise runs had error +2.8 +/- 1.1 mm and counterclockwise runs had
-3.4 +/- 1.6 mm. 14 of 16 runs converged to within the 5 mm band before the halfway
point; the two non-converging runs are excluded from the aggregate statistics. The
uncompensated controller was better in the first 30 s but drifted afterwards; this
matches the thermal time constant of the sensor mount.

Testing in a corroded pipe section would require access to the decommissioned line,
which was not available during the test campaign.

\begin{figure}
\caption{The wall-distance error for both controllers. See text for details.}
\end{figure}
